% Options for packages loaded elsewhere
\PassOptionsToPackage{unicode}{hyperref}
\PassOptionsToPackage{hyphens}{url}
\documentclass[
  11pt,
]{article}
\usepackage{xcolor}
\usepackage[margin=18mm]{geometry}
\usepackage{amsmath,amssymb}
\usepackage{graphicx}
\usepackage{tikz}
\usetikzlibrary{arrows.meta,positioning,calc}
\setcounter{secnumdepth}{-\maxdimen} % remove section numbering
\usepackage{iftex}
\ifPDFTeX
  \usepackage[T2A]{fontenc}
  \usepackage[utf8]{inputenc}
  \usepackage{textcomp} % provide euro and other symbols
\else % if luatex or xetex
  \usepackage{unicode-math} % this also loads fontspec
  \defaultfontfeatures{Scale=MatchLowercase}
  \defaultfontfeatures[\rmfamily]{Ligatures=TeX,Scale=1}
\fi
\usepackage{lmodern}
\ifPDFTeX\else
  % xetex/luatex font selection
\fi
% Use upquote if available, for straight quotes in verbatim environments
\IfFileExists{upquote.sty}{\usepackage{upquote}}{}
\IfFileExists{microtype.sty}{% use microtype if available
  \usepackage[]{microtype}
  \UseMicrotypeSet[protrusion]{basicmath} % disable protrusion for tt fonts
}{}
\makeatletter
\@ifundefined{KOMAClassName}{% if non-KOMA class
  \IfFileExists{parskip.sty}{%
    \usepackage{parskip}
  }{% else
    \setlength{\parindent}{0pt}
    \setlength{\parskip}{6pt plus 2pt minus 1pt}}
}{% if KOMA class
  \KOMAoptions{parskip=half}}
\makeatother
\usepackage{color}
\usepackage{fancyvrb}
\newcommand{\VerbBar}{|}
\newcommand{\VERB}{\Verb[commandchars=\\\{\}]}
\DefineVerbatimEnvironment{Highlighting}{Verbatim}{commandchars=\\\{\}}
% Add ',fontsize=\small' for more characters per line
\newenvironment{Shaded}{}{}
\newcommand{\AlertTok}[1]{\textcolor[rgb]{1.00,0.00,0.00}{\textbf{#1}}}
\newcommand{\AnnotationTok}[1]{\textcolor[rgb]{0.38,0.63,0.69}{\textbf{\textit{#1}}}}
\newcommand{\AttributeTok}[1]{\textcolor[rgb]{0.49,0.56,0.16}{#1}}
\newcommand{\BaseNTok}[1]{\textcolor[rgb]{0.25,0.63,0.44}{#1}}
\newcommand{\BuiltInTok}[1]{\textcolor[rgb]{0.00,0.50,0.00}{#1}}
\newcommand{\CharTok}[1]{\textcolor[rgb]{0.25,0.44,0.63}{#1}}
\newcommand{\CommentTok}[1]{\textcolor[rgb]{0.38,0.63,0.69}{\textit{#1}}}
\newcommand{\CommentVarTok}[1]{\textcolor[rgb]{0.38,0.63,0.69}{\textbf{\textit{#1}}}}
\newcommand{\ConstantTok}[1]{\textcolor[rgb]{0.53,0.00,0.00}{#1}}
\newcommand{\ControlFlowTok}[1]{\textcolor[rgb]{0.00,0.44,0.13}{\textbf{#1}}}
\newcommand{\DataTypeTok}[1]{\textcolor[rgb]{0.56,0.13,0.00}{#1}}
\newcommand{\DecValTok}[1]{\textcolor[rgb]{0.25,0.63,0.44}{#1}}
\newcommand{\DocumentationTok}[1]{\textcolor[rgb]{0.73,0.13,0.13}{\textit{#1}}}
\newcommand{\ErrorTok}[1]{\textcolor[rgb]{1.00,0.00,0.00}{\textbf{#1}}}
\newcommand{\ExtensionTok}[1]{#1}
\newcommand{\FloatTok}[1]{\textcolor[rgb]{0.25,0.63,0.44}{#1}}
\newcommand{\FunctionTok}[1]{\textcolor[rgb]{0.02,0.16,0.49}{#1}}
\newcommand{\ImportTok}[1]{\textcolor[rgb]{0.00,0.50,0.00}{\textbf{#1}}}
\newcommand{\InformationTok}[1]{\textcolor[rgb]{0.38,0.63,0.69}{\textbf{\textit{#1}}}}
\newcommand{\KeywordTok}[1]{\textcolor[rgb]{0.00,0.44,0.13}{\textbf{#1}}}
\newcommand{\NormalTok}[1]{#1}
\newcommand{\OperatorTok}[1]{\textcolor[rgb]{0.40,0.40,0.40}{#1}}
\newcommand{\OtherTok}[1]{\textcolor[rgb]{0.00,0.44,0.13}{#1}}
\newcommand{\PreprocessorTok}[1]{\textcolor[rgb]{0.74,0.48,0.00}{#1}}
\newcommand{\RegionMarkerTok}[1]{#1}
\newcommand{\SpecialCharTok}[1]{\textcolor[rgb]{0.25,0.44,0.63}{#1}}
\newcommand{\SpecialStringTok}[1]{\textcolor[rgb]{0.73,0.40,0.53}{#1}}
\newcommand{\StringTok}[1]{\textcolor[rgb]{0.25,0.44,0.63}{#1}}
\newcommand{\VariableTok}[1]{\textcolor[rgb]{0.10,0.09,0.49}{#1}}
\newcommand{\VerbatimStringTok}[1]{\textcolor[rgb]{0.25,0.44,0.63}{#1}}
\newcommand{\WarningTok}[1]{\textcolor[rgb]{0.38,0.63,0.69}{\textbf{\textit{#1}}}}
\usepackage{longtable,booktabs,array}
\newcounter{none} % for unnumbered tables
\usepackage{calc} % for calculating minipage widths
% Correct order of tables after \paragraph or \subparagraph
\usepackage{etoolbox}
\makeatletter
\patchcmd\longtable{\par}{\if@noskipsec\mbox{}\fi\par}{}{}
\makeatother
% Allow footnotes in longtable head/foot
\IfFileExists{footnotehyper.sty}{\usepackage{footnotehyper}}{\usepackage{footnote}}
\makesavenoteenv{longtable}
\ifLuaTeX
  \usepackage{luacolor}
  \usepackage[soul]{lua-ul}
\else
  \usepackage{soul}
\fi
\ifLuaTeX
\usepackage[bidi=basic,shorthands=off]{babel}
\else
\usepackage[bidi=default,shorthands=off]{babel}
\fi
\ifLuaTeX
  \usepackage{selnolig} % disable illegal ligatures
\fi
\setlength{\emergencystretch}{3em} % prevent overfull lines
\providecommand{\tightlist}{%
  \setlength{\itemsep}{0pt}\setlength{\parskip}{0pt}}
\usepackage{bookmark}
\IfFileExists{xurl.sty}{\usepackage{xurl}}{} % add URL line breaks if available
\urlstyle{same}
\hypersetup{
  pdftitle={Чорновик до залікової контрольної роботи},
  pdflang={ukrainian},
  hidelinks,
  pdfcreator={LaTeX via pandoc}}

\title{Чорновик до залікової контрольної роботи}
\author{}
\date{}

\begin{document}
\maketitle

\begin{quote}
Робоча чернетка. Деякі місця залишені як виправлення, щоб було видно хід думки, а не чистовий звіт.
\end{quote}

\subsection{Теоретична частина}\label{ux442ux435ux43eux440ux435ux442ux438ux447ux43dux430-ux447ux430ux441ux442ux438ux43dux430}

\begin{enumerate}
\def\labelenumi{\arabic{enumi}.}
\item
  Перевести двійкове число \texttt{1000\ 1011\ 1110\ 0011} у десяткову систему:

  Спочатку групую по 4 біти:

  \texttt{1000\ 1011\ 1110\ 0011\_2\ =\ 8BE3\_16}.

  Далі:

  \texttt{8BE3\_16\ =\ 8*16\^{}3\ +\ 11*16\^{}2\ +\ 14*16\ +\ 3\ =\ 35811\_10}.

  Відповідь: \texttt{35811}.
\item
  Перевести двійкове число \texttt{1000\ 1011\ 1110\ 0011} у шістнадцяткову систему:

  \texttt{1000\ 1011\ 1110\ 0011\_2\ =\ 8BE3\_16}.
\item
  За часовими діаграмами:

  \texttt{x\ =\ І-НЕ}, \texttt{y\ =\ АБО-НЕ}.

  \st{Спочатку подумав, що y може бути АБО}, але за діаграмою видно інверсію, тому \texttt{y\ =\ АБО-НЕ}.
\item
  Ознаки, що формуються в арифметико-логічних блоках під час арифметичних операцій:

  Правильний варіант: \texttt{Z\ =\ 1} визначається операцією \texttt{АБО-НЕ} над усіма розрядами результату; знак результату \texttt{N} дорівнює значенню старшого знакового розряду результату в модифікованому машинному коді \texttt{ЗР(n+2)}.
\item
  Ознака переповнення в знаковому розряді:

  Вибрати:

  \begin{itemize}
  \tightlist
  \item
    \texttt{flag\{Overflow\}\ =\ C\_\{n-1\}\ XOR\ C\_n}, де \texttt{C\_\{n-1\}} - перенос у знаковий розряд, \texttt{C\_n} - перенос із знакового розряду за межі розрядної сітки.
  \item
    Ознака переповнення встановлюється в \texttt{1}, коли старший знаковий біт обох доданків однаковий, але знак результату відрізняється.
  \end{itemize}
\item
  Навіщо в АЛП операції з дійсними числами виконуються над нормалізованими мантисами:

  Для забезпечення однакового формату чисел і точності обчислень, а також для уникнення втрати значущих розрядів при додаванні або відніманні.
\item
  Нормалізована мантиса двійкового числа:

  Вибрати:

  \begin{itemize}
  \tightlist
  \item
    Це дробова частина числа з плаваючою комою, яка містить його значущі біти.
  \item
    Нормалізовані мантиси додатних чисел у розряді з вагою \texttt{2\^{}-1} завжди мають одиницю.
  \end{itemize}
\item
  Чому неможливо виконувати алгебраїчне підсумовування і віднімання на звичайних суматорах у прямих машинних кодах:

  Правильний варіант: у прямому коді знакові розряди і основні розряди мають оброблятися по-різному. Віднімання на звичайному суматорі виконують, якщо аргументи подані у зворотному або доповняному коді; тоді віднімання замінюється додаванням числа з протилежним знаком.
\item
  Додавання двійкових мантис зі знаками:

  Дано:

  \texttt{A\ =\ 0\ 10101}, \texttt{B\ =\ 0\ 0011010}.

  Обидва числа додатні, тому додаємо мантиси:

\begin{Shaded}
\begin{Highlighting}[]
\NormalTok{  0.1010100}
\NormalTok{+ 0.0011010}
\NormalTok{{-}{-}{-}{-}{-}{-}{-}{-}{-}{-}{-}}
\NormalTok{  0.1101110}
\end{Highlighting}
\end{Shaded}

  Точний результат: \texttt{0\ 1101110}.

  У 6-розрядній сітці без округлення: \texttt{0\ 11011}.

  \st{Можна було б округлити}, але в умові просто записати в 6-розрядну сітку, тому беру обрізання.
\item
  Яка з мікросхем має пам'ять:

  \texttt{Лічильник}.
\item
  Чим відрізняється однорозрядний напівсуматор від повного однорозрядного суматора:

  Однорозрядний суматор за розрядним значенням \texttt{xi} і \texttt{yi} доданків і за значенням переносу \texttt{zi} із сусіднього молодшого розряду формує значення розрядної суми \texttt{Si} і переносу у наступний старший розряд \texttt{Pi}. Однорозрядний напівсуматор не враховує вхідний перенос \texttt{zi}.
\item
  Восьмирозрядний суматор складається із:

  \texttt{Ланцюга\ із\ восьми\ повних\ однобітних\ суматорів}.
\item
  Дешифратор / шифратор: правдиві твердження:

  Вибрати:

  \begin{itemize}
  \tightlist
  \item
    Неповний дешифратор має менше ніж \texttt{2\^{}n} виходів.
  \item
    Шифратор - це пристрій, що виконує функцію, зворотну дешифратору.
  \item
    Повний дешифратор має \texttt{n} входів і \texttt{2\^{}n} виходів.
  \item
    При подачі сигналу на один із входів шифратора, тобто унітарного коду, на виході повинен утворитися двійковий код.
  \item
    Дешифратор - це комбінаційна схема для перетворення двійкового позиційного коду на унітарний.
  \end{itemize}
\item
  Мікрооперації, які можуть виконуватися на регістрах:

  Вибрати:

  \begin{itemize}
  \tightlist
  \item
    Зсув логічний.
  \item
    Зсув арифметичний.
  \item
    Зберігання.
  \item
    Запис.
  \item
    Встановлення в нуль, тобто очищення.
  \end{itemize}
\end{enumerate}

\subsection{Функціональне моделювання в САПР ModelSim}\label{ux444ux443ux43dux43aux446ux456ux43eux43dux430ux43bux44cux43dux435-ux43cux43eux434ux435ux43bux44eux432ux430ux43dux43dux44f-ux432-ux441ux430ux43fux440-modelsim}

\begin{enumerate}
\def\labelenumi{\arabic{enumi}.}
\setcounter{enumi}{14}
\item
  Як можна реалізувати 16-розрядний регістр на ПЛІС (FPGA):

  Знадобиться 16 логічних блоків ПЛІС, де кожен біт регістра реалізується за допомогою одного \texttt{D}-тригера (\texttt{flip-flop}).
\item
  З якою ціллю розробляється Testbench у САПР:

  Для емуляції оточення розроблюваного пристрою і аналізу результатів.
\item
  Яку послідовність вхідних сигналів сформує Testbench:

  Код:

\begin{Shaded}
\begin{Highlighting}[]
\OperatorTok{@(}\KeywordTok{posedge}\NormalTok{ clk}\OperatorTok{);}\NormalTok{ out1}\OperatorTok{=}\DecValTok{0}\OperatorTok{;}\NormalTok{ out2}\OperatorTok{=}\DecValTok{0}\OperatorTok{;}
\OperatorTok{@(}\KeywordTok{posedge}\NormalTok{ clk}\OperatorTok{);}\NormalTok{ out1}\OperatorTok{=}\DecValTok{0}\OperatorTok{;}\NormalTok{ out2}\OperatorTok{=}\DecValTok{1}\OperatorTok{;}
\OperatorTok{@(}\KeywordTok{posedge}\NormalTok{ clk}\OperatorTok{);}\NormalTok{ out1}\OperatorTok{=}\DecValTok{1}\OperatorTok{;}\NormalTok{ out2}\OperatorTok{=}\DecValTok{0}\OperatorTok{;}
\OperatorTok{@(}\KeywordTok{posedge}\NormalTok{ clk}\OperatorTok{);}\NormalTok{ out1}\OperatorTok{=}\DecValTok{1}\OperatorTok{;}\NormalTok{ out2}\OperatorTok{=}\DecValTok{1}\OperatorTok{;}
\end{Highlighting}
\end{Shaded}

  Послідовність:

  \texttt{00\ 01\ 10\ 11}.
\item
  Таблиця істинності у часовій області для неповного дешифратора 3-to-8 з входом \texttt{EN}:

  Правильний варіант: \texttt{2}.

  Логіка така: при \texttt{EN\ =\ 0} вихід \texttt{Y\ =\ 00000000}, при \texttt{EN\ =\ 1} активується один вихід за адресою \texttt{A2\ A1\ A0}.
\item
  Скільки тактів буде тривати процес моделювання заданим Testbench:

  Правильний варіант: \texttt{вісім\ тактів}.

  Пояснення: 4 такти для комбінацій \texttt{00}, \texttt{01}, \texttt{10}, \texttt{11}, потім ще 1 такт окремо і \texttt{repeat(3)} додає ще 3 такти.

  \texttt{4\ +\ 1\ +\ 3\ =\ 8}.
\item
  Для тестування якого пристрою використовується Testbench з \texttt{reg\ {[}1:0{]}\ sel}, \texttt{reg\ {[}3:0{]}\ d}, \texttt{wire\ y}:

  Правильний варіант: \texttt{мультиплексор\ 4x1}.

  Пояснення: є 4 інформаційні входи \texttt{d{[}3:0{]}}, 2-бітний вибір \texttt{sel} і один вихід \texttt{y}.
\end{enumerate}

\subsection{Практична частина}\label{ux43fux440ux430ux43aux442ux438ux447ux43dux430-ux447ux430ux441ux442ux438ux43dux430}

\begin{enumerate}
\def\labelenumi{\arabic{enumi}.}
\setcounter{enumi}{20}
\item
  Адресний селектор зовнішнього пристрою:

  Правильний варіант: \texttt{PC\ =\ 29A4h}.

  Пояснення: інверсні входи селектора мають отримати \texttt{0}, прямі входи - \texttt{1}; \texttt{A0} не впливає на вибір, за умовою беремо \texttt{A0\ =\ 0}.
\item
  Практичне завдання за номером студентського квитка:

  Номер: \texttt{4106}.

  Переведення у двійковий код:

\begin{Shaded}
\begin{Highlighting}[]
\NormalTok{4106\_10 = 1000000001010\_2}
\end{Highlighting}
\end{Shaded}

  Далі вирівнюю число праворуч у 16-розрядній сітці, тобто дописую нулі ліворуч:

\begin{Shaded}
\begin{Highlighting}[]
\NormalTok{4106\_10 = 0001 0000 0000 1010\_2}
\end{Highlighting}
\end{Shaded}

  Два 8-розрядні аргументи:

\begin{Shaded}
\begin{Highlighting}[]
\NormalTok{R1 = 00010000\_2 = +16}
\NormalTok{R2 = 00001010\_2 = +10}
\end{Highlighting}
\end{Shaded}

  Обидва аргументи додатні, у доповняльному коді запис не змінюється:

\begin{Shaded}
\begin{Highlighting}[]
\NormalTok{A = 00010000}
\NormalTok{B = 00001010}
\end{Highlighting}
\end{Shaded}

  Виконується додавання:

\begin{Shaded}
\begin{Highlighting}[]
\NormalTok{  00010000}
\NormalTok{+ 00001010}
\NormalTok{{-}{-}{-}{-}{-}{-}{-}{-}{-}{-}}
\NormalTok{  00011010 = 26}
\end{Highlighting}
\end{Shaded}

  У 9-розрядному регістрі:

\begin{Shaded}
\begin{Highlighting}[]
\NormalTok{S = 000011010}
\end{Highlighting}
\end{Shaded}

  Переповнення немає, бо \texttt{A{[}7{]}\ =\ 0}, \texttt{B{[}7{]}\ =\ 0}, \texttt{S{[}7{]}\ =\ 0}.

  Далі виконується гілка без переповнення:

\begin{Shaded}
\begin{Highlighting}[]
\NormalTok{S }\OperatorTok{=}\NormalTok{ S }\OperatorTok{\textgreater{}\textgreater{}\textgreater{}} \DecValTok{2}\OperatorTok{;}
\NormalTok{RESULT }\OperatorTok{=}\NormalTok{ S}\OperatorTok{[}\DecValTok{8}\OperatorTok{:}\DecValTok{7}\OperatorTok{];}
\end{Highlighting}
\end{Shaded}

  Після арифметичного зсуву праворуч на 2:

\begin{Shaded}
\begin{Highlighting}[]
\NormalTok{до зсуву:    S = 000011010 = 26}
\NormalTok{після >>> 1: S = 000001101 = 13}
\NormalTok{S = 000000110 = 6}
\NormalTok{RESULT = S[8:7] = 00}
\end{Highlighting}
\end{Shaded}

  Цифрова діаграма станів регістрів:

  {\def\LTcaptype{none}
  \begin{longtable}[]{@{}rlccccc@{}}
  \toprule\noalign{}
  Етап & Операція & A & B & S & OVERFLOW & RESULT \\
  \midrule\noalign{}
  \endhead
  \bottomrule\noalign{}
  \endlastfoot
  0 & Початкові аргументи & \texttt{00010000} & \texttt{00001010} & -- & -- & -- \\
  1 & \texttt{S = A + B} & \texttt{00010000} & \texttt{00001010} & \texttt{000011010} & -- & -- \\
  2 & Перевірка overflow & \texttt{00010000} & \texttt{00001010} & \texttt{000011010} & \texttt{0} & -- \\
  3 & \texttt{S = S >>> 2} & \texttt{00010000} & \texttt{00001010} & \texttt{000000110} & \texttt{0} & -- \\
  4 & \texttt{RESULT = S[8:7]} & \texttt{00010000} & \texttt{00001010} & \texttt{000000110} & \texttt{0} & \texttt{00} \\
  \end{longtable}
  }

  Регістр \texttt{C} у цьому варіанті не змінюється, бо гілка з \texttt{C\ =\ S} виконується тільки при переповненні.

  Ознаки для \texttt{S}:

\begin{Shaded}
\begin{Highlighting}[]
\NormalTok{S = 000000110}
\NormalTok{zero\_flag     = 0, бо S не дорівнює нулю}
\NormalTok{overflow\_flag = 0, бо переповнення не було}
\NormalTok{sign\_flag     = 0, бо знаковий біт дорівнює 0}
\end{Highlighting}
\end{Shaded}

  Результат у прямому коді: \texttt{00000110\_2\ =\ +6}.

  Перевірка десятковими еквівалентами:

\begin{Shaded}
\begin{Highlighting}[]
\NormalTok{R1 = +16}
\NormalTok{R2 = +10}
\NormalTok{S  = 16 + 10 = 26}
\NormalTok{S >>> 2 = 26 / 4 = 6}
\end{Highlighting}
\end{Shaded}

  У Google Forms для \texttt{RESULT}: \texttt{00}.
\item
  Сторінкова пам'ять:

  Умова: загальний об'єм пам'яті \texttt{516K\ x\ 16}, об'єм однієї сторінки пам'яті \texttt{128K\ x\ 16}.

  Спочатку визначаю шину даних. Другий множник у записі пам'яті означає розрядність одного слова:

\begin{Shaded}
\begin{Highlighting}[]
\NormalTok{n = 16}
\end{Highlighting}
\end{Shaded}

  Тобто \texttt{ШД\ =\ 16\ біт}.

  Одна сторінка має об'єм \texttt{128K\ x\ 16}. Кількість адресованих слів на сторінці:

\begin{Shaded}
\begin{Highlighting}[]
\NormalTok{128K = 128 * 1024 = 2\^{}7 * 2\^{}10 = 2\^{}17}
\NormalTok{m2 = 17}
\end{Highlighting}
\end{Shaded}

  Отже молодші біти адреси \texttt{A16..A0} задають адресу всередині сторінки.

  Загальний об'єм пам'яті:

\begin{Shaded}
\begin{Highlighting}[]
\NormalTok{516K = 516 * 1024 = 528384}
\NormalTok{2\^{}19 = 524288 \textless{} 528384 \textless{} 2\^{}20}
\end{Highlighting}
\end{Shaded}

  Тому для повної адресації потрібно 20 адресних бітів:

\begin{Shaded}
\begin{Highlighting}[]
\NormalTok{m = 20}
\end{Highlighting}
\end{Shaded}

  Тобто \texttt{ША\ =\ 20\ біт}.

  Кількість сторінок:

\begin{Shaded}
\begin{Highlighting}[]
\NormalTok{N = ceil(516K / 128K)}
\NormalTok{N = ceil(4.03125)}
\NormalTok{N = 5 сторінок}
\end{Highlighting}
\end{Shaded}

  Тут саме 5 сторінок, бо 4 сторінки дають тільки:

\begin{Shaded}
\begin{Highlighting}[]
\NormalTok{4 * 128K = 512K}
\end{Highlighting}
\end{Shaded}

  а потрібно \texttt{516K}. \st{Спочатку можна було б взяти 4 сторінки}, але це покриває тільки \texttt{512K}, тому потрібна ще частина п'ятої сторінки.

  Для вибору 5 сторінок треба дешифратор із трьома входами:

\begin{Shaded}
\begin{Highlighting}[]
\NormalTok{m1 = ceil(log2(5)) = 3}
\end{Highlighting}
\end{Shaded}

  Старші адресні біти для вибору сторінки: \texttt{A19,\ A18,\ A17}.

  Підсумкова таблиця параметрів:

  {\def\LTcaptype{none}
  \begin{longtable}[]{@{}lr@{}}
  \toprule\noalign{}
  Параметр & Значення \\
  \midrule\noalign{}
  \endhead
  \bottomrule\noalign{}
  \endlastfoot
  Загальний об'єм пам'яті & \texttt{516K\ x\ 16} \\
  Об'єм однієї сторінки & \texttt{128K\ x\ 16} \\
  Розрядність шини даних, ШД & \texttt{16} \\
  Розрядність шини адреси, ША & \texttt{20} \\
  Кількість сторінок \(N\) & \texttt{5} \\
  Кількість входів дешифратора \(m_1\) & \texttt{3} \\
  Розрядність адреси сторінки \(m_2\) & \texttt{17} \\
  \end{longtable}
  }

  Таблиця істинності дешифратора \texttt{3\ -\textgreater{}\ 8} для 5 сторінок:

  {\def\LTcaptype{none} % do not increment counter
  \begin{longtable}[]{@{}rrrrrrrrl@{}}
  \toprule\noalign{}
  A19 & A18 & A17 & CS0 & CS1 & CS2 & CS3 & CS4 & Сторінка \\
  \midrule\noalign{}
  \endhead
  \bottomrule\noalign{}
  \endlastfoot
  0 & 0 & 0 & 1 & 0 & 0 & 0 & 0 & 0 \\
  0 & 0 & 1 & 0 & 1 & 0 & 0 & 0 & 1 \\
  0 & 1 & 0 & 0 & 0 & 1 & 0 & 0 & 2 \\
  0 & 1 & 1 & 0 & 0 & 0 & 1 & 0 & 3 \\
  1 & 0 & 0 & 0 & 0 & 0 & 0 & 1 & 4 \\
  1 & 0 & 1 & 0 & 0 & 0 & 0 & 0 & не використ. \\
  1 & 1 & 0 & 0 & 0 & 0 & 0 & 0 & не використ. \\
  1 & 1 & 1 & 0 & 0 & 0 & 0 & 0 & не використ. \\
  \end{longtable}
  }

  Логічні рівняння:

\begin{Shaded}
\begin{Highlighting}[]
\NormalTok{CS0 = \textasciitilde{}A19 \& \textasciitilde{}A18 \& \textasciitilde{}A17}
\NormalTok{CS1 = \textasciitilde{}A19 \& \textasciitilde{}A18 \&  A17}
\NormalTok{CS2 = \textasciitilde{}A19 \&  A18 \& \textasciitilde{}A17}
\NormalTok{CS3 = \textasciitilde{}A19 \&  A18 \&  A17}
\NormalTok{CS4 =  A19 \& \textasciitilde{}A18 \& \textasciitilde{}A17}
\end{Highlighting}
  \end{Shaded}

  Структурна схема багатосторінкової пам'яті:

  \begin{center}
  \resizebox{\linewidth}{!}{%
  \begin{tikzpicture}[
      font=\small,
      block/.style={draw, rectangle, align=center, minimum width=3.4cm, minimum height=1.35cm},
      mem/.style={draw, rectangle, align=center, minimum width=2.25cm, minimum height=1.15cm},
      bus/.style={-Latex, very thick},
      sig/.style={-Latex, thin},
      tap/.style={thin}
  ]
      \node[block, minimum height=1.7cm] (dec) at (0,2.2)
          {Дешифратор\\\(3 \rightarrow 8\)\\використано 5 виходів};

      \node[left=3.6cm of dec] (addrhi) {\(A_{19}, A_{18}, A_{17}\)};
      \draw[bus] (addrhi.east) -- node[above]{старші біти адреси} (dec.west);

      \foreach \i/\x in {0/-5.0,1/-2.5,2/0,3/2.5,4/5.0} {
          \node[mem] (p\i) at (\x,-1.0) {Сторінка \(\i\)\\\(128K \times 16\)};
      }

      \draw[tap] (dec.south) -- (0,0.95);
      \foreach \i/\x in {0/-5.0,1/-2.5,2/0,3/2.5,4/5.0} {
          \draw[sig] (0,0.95) -| node[pos=0.23, above]{\(CS_{\i}\)} (p\i.north);
      }

      \draw[bus] (-6.35,-2.65) -- node[below]{молодші біти адреси \(A_{16..0}\), 17 біт} (6.35,-2.65);
      \foreach \i/\x in {0/-5.0,1/-2.5,2/0,3/2.5,4/5.0} {
          \draw[tap] (\x,-2.65) -- (\x,-1.58);
      }

      \draw[bus] (-6.35,-3.45) -- node[below]{спільна шина даних \(D_{15..0}\), 16 біт} (6.35,-3.45);
      \foreach \i/\x in {0/-5.0,1/-2.5,2/0,3/2.5,4/5.0} {
          \draw[tap] (\x+0.42,-3.45) -- (\x+0.42,-1.58);
      }

      \node[align=center] at (0,3.55)
          {Старші біти вибирають сторінку, молодші біти адресують слово всередині сторінки};
  \end{tikzpicture}}
  \end{center}

  Пояснення підключення: старші біти адреси \texttt{A19..A17} подаються на дешифратор вибору сторінки, молодші \texttt{A16..A0} йдуть паралельно на всі сторінки пам'яті, а спільна шина даних \texttt{D15..D0} підключена до виходів сторінок. Активний сигнал \texttt{CSi} дозволяє роботу тільки однієї сторінки. Комбінації \texttt{101}, \texttt{110}, \texttt{111} не використовуються.

  У Google Forms:

\begin{Shaded}
\begin{Highlighting}[]
\NormalTok{Кількість сторінок: 5}
\NormalTok{Розрядність даних ШД: 16}
\NormalTok{Розрядність адреси ША: 20}
\NormalTok{Кількість входів дешифратора m1: 3}
\NormalTok{Розрядність адреси сторінки m2: 17}
\end{Highlighting}
\end{Shaded}
\end{enumerate}

\end{document}
